\section{Qué se perdió al atravesar el material}
La pérdida de inserción describe la reducción de potencia que llega al otro lado de una lámina. El caso de agua marina de 7 mm permite distinguir sus causas. A 5 px/mm se reflejó el 62.78\,\% de la potencia incidente, se absorbió el 35.20\,\% y se transmitió el 1.92\,\%. En músculo, al mismo espesor y resolución, las fracciones fueron 31.64\,\%, 53.21\,\% y 15.03\,\%. Aunque el agua presentó la mayor atenuación por unidad de distancia, el músculo absorbió una fracción mayor de la potencia incidente en esta configuración. La reflexión en la primera interfaz condicionó cuánta potencia pudo entrar a cada lámina. El cálculo analítico incluyó también las reflexiones internas.

Esta interpretación depende de la geometría. En la antena inmersa, la energía que entró en el anillo ya se propagaba dentro del medio. En la lámina, la potencia de referencia incidió desde el vacío. Por ello, el 93.75\,\% de disipación local del monopolo en agua y el 35.20\,\% de absorción de la lámina responden a preguntas físicas distintas.

\section{Qué explicó la comparación con la teoría}
El control de onda plana recuperó la atenuación esperada con una discrepancia máxima del 0.146\,\% a 5 px/mm en los medios con pérdidas. En el monopolo, la mayor discrepancia alcanzó el 5.18\,\% en músculo. En este segundo problema se extrajo una pendiente sobre una línea radial finita. La multiplicación por $r^2$ descontó la expansión esperada en la zona de radiación, pero no eliminó toda la estructura del campo de una antena finita.

El vacío permitió apreciar ese límite. Con conductividad nula, la integral de Joule se anuló. Sin embargo, la pendiente aparente del monopolo fue ligeramente negativa y su intervalo condicional excluyó cero. La onda plana en vacío sí dio un intervalo compatible con cero. Por consiguiente, el pequeño crecimiento del perfil compensado correspondió a la distribución espacial observada y al procedimiento de extracción. Los datos no aislaron una causa única entre geometría, ventana y discretización.

\section{Precisión del ajuste y sensibilidad de malla}
En agua a 5 px/mm, el IC del 95\,\% pasó de [1.18615, 1.18705] a [1.18595, 1.18719] dB/mm al cuadruplicar el bloque. En esa misma malla, variar el inicio de la ventana llevó la pendiente de 1.18520 a 1.18872 dB/mm. El intervalo por bloques describió la variación de la estimación bajo remuestreo del perfil. La diferencia con la onda plana y los cambios entre mallas superaron la anchura de ese intervalo y se evaluaron por separado. La variabilidad de las propiedades entre muestras físicas quedó sin cuantificar.

La pendiente en agua cambió en un solo sentido en las cuatro mallas. Sin embargo, los incrementos entre 5 y 10 y entre 10 y 20 px/mm fueron casi iguales. El cociente de esas diferencias produjo un orden aparente de apenas 0.018. En estas condiciones la extrapolación depende demasiado de pequeñas variaciones en las pendientes para asignar un error absoluto al continuo. El valor a 20 px/mm se utilizó como referencia numérica. En músculo y concreto, dos mallas permitieron calcular sensibilidad entre resoluciones.

\section{Alcance de la resolución común}
A 5 px/mm se obtuvieron aproximadamente 58 celdas por longitud de onda en agua, 69 en músculo, 218 en concreto y 500 en vacío. El diámetro del conductor ocupó solo dos celdas. Al comparar con la malla más fina de cada material, la pendiente presentó diferencias relativas de 0.81\,\% en agua, 2.39\,\% en músculo y 0.025\,\% en concreto. Las diferencias relativas en la fracción disipada fueron 0.099\,\%, 0.404\,\% y 0.025\,\%. Estos valores sustentaron la comparación principal a 5 px/mm para ambas magnitudes.

Las potencias entrantes a esa resolución superaron a las de la malla más fina en 23.9\,\% en agua, 75.2\,\% en músculo y 50.9\,\% en concreto. Un balance puede cerrar con precisión en cada malla aunque su entrada y su salida cambien juntas al refinarla. El pequeño diámetro discretizado del conductor y la representación de la alimentación pudieron influir en esos cambios. Los controles de espesor absorbente produjeron variaciones de pendiente inferiores a $10^{-5}$ dB/mm en los tres medios evaluados. Su efecto fue menor que el del refinamiento en estos casos.

La corrida en vacío a 5 px/mm requirió unas 5.49 horas y 4.76 GiB de memoria residente máxima. La de concreto a 10 px/mm requirió 8.68 horas y 4.17 GiB. La ejecución secuencial con un límite de 5 GiB condicionó el refinamiento disponible. En vacío se contó con una sola malla del monopolo, por lo que la sensibilidad espacial de esa geometría permaneció abierta.

\section{Evaluación de los objetivos}
La pérdida de inserción se predijo para capas homogéneas de espesor definido a incidencia normal. Los espesores reservados permitieron comprobar la formulación frente a la solución analítica sin reajustar parámetros. Se propuso un procedimiento reproducible que fija la referencia incidente, la frecuencia, la geometría, los monitores, la ventana radial y el criterio de estacionariedad. Los anexos permiten reconstruir los cálculos y comparar las mallas con las mismas definiciones.

Para un modelo de propagación, la pendiente del monopolo proporciona un término de pérdida radial relativa en dB/mm dentro de la ventana estudiada. Su traslado a un enlace de telecomunicaciones requiere conocer la potencia de la fuente, el patrón de radiación, las interfaces del entorno y la respuesta del receptor. La campaña estableció el alcance físico y numérico de ese término a 3 GHz. La validación frente a mediciones de enlaces queda como trabajo posterior.
